% ----------------------------------------------------------
% Introdução
% ----------------------------------------------------------
\chapter{Introdução}
\label{cap:introducao}
% ----------------------------------------------------------
A teoria da computação estuda as propriedades dos sistemas computacionais. Esta teoria foi formalizada pela primeira vez, de forma efetiva, em 1936, por Alan Turing e Alonzo Church, que construíram as bases do que hoje chamamos de computação, por meio da elaboração do conceito de linguagens formais.

Linguagens formais são lidas por meio de máquinas de estados e escritas por meio de gramáticas formais. Neste trabalho, serão enfatizadas as máquinas de estados, mais especificamente autômatos de pilha não determinísticos (APN), por serem particularmente mais abstratos e de representação mais complexa que as demais máquinas de estados.

Sendo a teoria da computação a base das tecnologias atuais, torna-se essencial que os discentes dos cursos da área de computação, como sistemas de informação e engenharia da computação, dominem esses conceitos para que se tornem bons profissionais. Uma vez que tais conceitos não apenas fornecem a base teórica para a computação, mas também promovem habilidades de pensamento crítico e análise rigorosa, conferindo a capacidade de formalizar problemas, definir algoritmos e analisar sua eficiência.

Como forma de auxiliar os discentes dos cursos da área da computação a visualizar de forma mais clara os autômatos de pilha não determinísticos, este trabalho tem como objetivo desenvolver uma ferramenta web para montagem e visualização de APNs em execução.


% ----------------------------

\section{Elaboração do capítulo}

Este capítulo apresenta o seu trabalho. Você deve contextualizar o problema abordado, descrever os objetivos gerais e específicos, apresentar a metodologia e como o trabalho está estruturado.

\section{O problema de pesquisa}
\label{sec:problema}

Ciente da importancia do estudo da teoria da computação, utilizada nos dias de hoje como ferramenta para o desenvolvimento dos compiladores das linguagens de programação modernas, é fundamental que os alunos dos cursos superiores da area de computação dominem esta diciplina. Contudo, eles normalmente precisam lidar com alguns desafios durante o estudo, sendo dois deles:

\begin{itemize}
	\item O estudo das maquinas de estados é realizado utilizando apenas diagramas em um papel e abstrair o seu funcionamento dada essa representação limitada pode se tornar um desafio;
	\item Criar uma maquina de estado que reconheça uma linguagem especifica e apenas ela, é como responder uma charada ou enigma, em que uma resposta precisa satisfazer uma serie de condições especificas para ser considerada correta.
\end{itemize}

Este trabalho visa amenizar o primeiro desafio e auxiliar no enfrentamento do segundo, facilitando o exercício pratico da disciplina.

\section{Objetivos}
\label{sec:objetivos}

O presente trabalho consiste em facilitar o aprendizado da disciplina de fundamentos teóricos da computação.

Este trabalho possui aos seguintes objetivos específicos:

\begin{itemize}
	\item Criar um software gráfico que auxilie os alunos simulando o funcionamento de maquinas de estado em tempo real;
	\item Disponibilizar esse sistema para alunos da disciplina de fundamentos teóricos da computação do campus ICEA;
	\item Avaliar o impacto do sistema no aprendizado dos alunos por meio de um questionário. 
\end{itemize}

\section{Metodologia}
\label{sec:metodologia}

O objeto de pesquisa deste trabalho ...

Os passos para execução deste trabalho são assim definidos:

\begin{itemize}
	\item Revisão da literatura
	\item Desenvolvimento do software de auxilio aos alunos
	\item Teste por parte dos alunos
	\item Análise do feedback dos alunos
\end{itemize}

\section{Organização do trabalho}

\textbf{É importante observar que a estrutura é apresentada a partir do próximo capítulo. O capítulo de Introdução não deve compor esta descrição. Além disso, sempre que você fizer referência à algum item específico, a inicial deve ser maiúscula. Por exemplo, Capítulo 2, Tabela 5, Figura 1, dentre outros.}

O restante deste trabalho é organizado como se segue. O Capítulo~\ref{cap:revisao} apresenta...
